\documentclass[11pt,a4paper]{article}

% ── Encoding ──────────────────────────────────────────────────────────────────
\usepackage[T1]{fontenc}
\usepackage[utf8]{inputenc}

% ── Page Layout ───────────────────────────────────────────────────────────────
\usepackage[a4paper, top=1in, bottom=1in, left=1in, right=1in]{geometry}

% ── Graphics ──────────────────────────────────────────────────────────────────
\usepackage{graphicx}
\usepackage{float}

% ── Tables ────────────────────────────────────────────────────────────────────
\usepackage{booktabs}
\usepackage{array}
\usepackage{tabularx}
\usepackage{longtable}
\usepackage{multirow}

% ── Mathematics ───────────────────────────────────────────────────────────────
\usepackage{amsmath}
\usepackage{amssymb}

% ── Colors ────────────────────────────────────────────────────────────────────
\usepackage[dvipsnames]{xcolor}
\definecolor{primary}{RGB}{15,23,42}
\definecolor{secondary}{RGB}{30,41,59}
\definecolor{accent}{RGB}{37,99,235}
\definecolor{lightbg}{RGB}{248,250,252}
\definecolor{codebg}{RGB}{240,244,248}
\definecolor{warnbg}{RGB}{255,243,205}
\definecolor{successbg}{RGB}{212,237,218}

% ── Hyperlinks ────────────────────────────────────────────────────────────────
\usepackage[hidelinks]{hyperref}
\hypersetup{
    colorlinks=true,
    linkcolor=accent,
    urlcolor=accent,
    citecolor=accent
}

% ── Code Listings ─────────────────────────────────────────────────────────────
\usepackage{listings}
\lstdefinestyle{HiveStyle}{
    language=SQL,
    basicstyle=\ttfamily\footnotesize,
    backgroundcolor=\color{codebg},
    commentstyle=\color{OliveGreen}\itshape,
    keywordstyle=\color{accent}\bfseries,
    stringstyle=\color{BrickRed},
    numberstyle=\tiny\color{gray},
    numbers=left,
    stepnumber=1,
    numbersep=8pt,
    breaklines=true,
    showstringspaces=false,
    frame=single,
    framesep=4pt,
    xleftmargin=12pt,
    rulecolor=\color{gray!40},
    tabsize=4,
    captionpos=b,
    morekeywords={
        CREATE, TABLE, DROP, IF, EXISTS, STORED, AS, ORC, TEXTFILE,
        TBLPROPERTIES, LOAD, DATA, INPATH, INTO, ROW, FORMAT, DELIMITED,
        FIELDS, TERMINATED, BY, SELECT, FROM, WHERE, GROUP, ORDER,
        HAVING, DISTINCT, PARTITION, OVER, LAG, ROW_NUMBER, ROUND,
        CAST, COALESCE, NULLIF, TRIM, LOWER, UPPER, INITCAP, SPLIT,
        REGEXP_REPLACE, CONCAT, POWER, WITH, COUNT, MAX, AVG, MIN,
        CASE, WHEN, THEN, ELSE, END, AND, OR, NOT, IN, NULL, INT,
        STRING, BIGINT, DOUBLE, ASC, DESC, LIMIT, SET, SNAPPY
    }
}
\lstset{style=HiveStyle}

% ── Headers / Footers ─────────────────────────────────────────────────────────
\usepackage{fancyhdr}
\setlength{\headheight}{15pt}
\pagestyle{fancy}
\fancyhf{}
\lhead{\small\color{secondary}Big Data Analytics – Group Assignment}
\rhead{\small\color{secondary}Hadoop \& Hive: Employee Dataset Analysis}
\cfoot{\thepage}
\renewcommand{\headrulewidth}{0.4pt}

% ── Section Formatting ────────────────────────────────────────────────────────
\usepackage{titlesec}
\titleformat{\section}{\Large\bfseries\color{primary}}{\thesection}{1em}{}[{\color{accent}\titlerule[0.6pt]}]
\titleformat{\subsection}{\large\bfseries\color{secondary}}{\thesubsection}{1em}{}
\titleformat{\subsubsection}{\normalsize\bfseries\color{secondary}}{\thesubsubsection}{1em}{}

% ── Misc ──────────────────────────────────────────────────────────────────────
\usepackage{parskip}
\usepackage{enumitem}
\usepackage{caption}
\captionsetup{font=small, labelfont=bf, labelsep=period}

% ── Document Metadata ─────────────────────────────────────────────────────────
\title{
    \vspace{-1cm}
    {\Huge\bfseries\color{primary} Big Data Analytics}\\[0.4cm]
    {\Large\color{secondary} Employee Dataset Analysis Using Apache Hadoop \& Hive}\\[0.3cm]
    {\normalsize\color{gray} Group Assignment Report}
}
\author{Big Data Analytics Group}
\date{\today}

% =============================================================================
\begin{document}
% =============================================================================

\maketitle
\thispagestyle{fancy}

\begin{abstract}
This report documents a full end-to-end Big Data analysis of an Employee Dataset
using Apache Hadoop (HDFS) and Apache Hive (HQL). The dataset of 100 employee
records is uploaded to HDFS, loaded into a Hive staging table, and processed
through a multi-stage pipeline. The analysis covers five data preprocessing
problems, job-title column splitting, PhD experience queries, skills-count
ordering, call-log growth rate computation, and call-log future value prediction
using a compound growth model.
\end{abstract}

\tableofcontents
\newpage

% ─────────────────────────────────────────────────────────────────────────────
\section{Introduction}
% ─────────────────────────────────────────────────────────────────────────────
Big Data analytics requires robust pipelines that can store, clean, and query
large volumes of structured data at scale. Apache Hadoop provides the distributed
storage (HDFS) and compute (MapReduce/YARN) foundation, while Apache Hive
translates SQL-like queries (HQL) into distributed jobs running on that cluster.

This practical demonstrates the full workflow on an \textbf{Employee Dataset}
(\texttt{EmployeeDataset.csv}, 100 rows, 12 columns) covering roles from AI
Engineer to Cybersecurity Analyst across countries, education levels, and
industries.

\subsection{Dataset Schema}
\begin{table}[H]
\centering
\caption{EmployeeDataset.csv -- Column Definitions}
\begin{tabularx}{\textwidth}{@{}llX@{}}
\toprule
\textbf{Column} & \textbf{Type} & \textbf{Description} \\
\midrule
\texttt{job\_id}          & INT    & Unique employee identifier \\
\texttt{job\_title}       & STRING & Full job title (e.g.\ \textit{AI Engineer}) \\
\texttt{experience\_years}& INT    & Years of professional experience \\
\texttt{education\_level} & STRING & Highest qualification (PhD, Bachelor, etc.) \\
\texttt{skills\_count}    & INT    & Number of listed technical skills \\
\texttt{industry}         & STRING & Sector (Healthcare, Finance, Technology…) \\
\texttt{company\_size}    & STRING & Small / Medium / Large / Enterprise / Startup \\
\texttt{location}         & STRING & Country of employment \\
\texttt{remote\_work}     & STRING & Yes / No / Hybrid \\
\texttt{call\_log}        & INT    & Number of logged calls \\
\texttt{salary}           & BIGINT & Annual salary (USD) \\
\texttt{date\_of\_marriage}& STRING & Marriage date (various formats) \\
\bottomrule
\end{tabularx}
\end{table}

% ─────────────────────────────────────────────────────────────────────────────
\section{System Architecture \& Environment Setup}
% ─────────────────────────────────────────────────────────────────────────────

\subsection{Pipeline Architecture}
The analysis pipeline follows four logical layers:
\begin{enumerate}[leftmargin=2cm]
    \item \textbf{Storage Layer (HDFS):} The raw CSV is uploaded to
          \texttt{/data/employee/} on HDFS.
    \item \textbf{Staging Layer (Hive):} A text-format external table
          \texttt{emp\_raw} is mapped over the HDFS path.
    \item \textbf{Processing Layer (HQL):} Window functions, CTAS queries,
          and analytical SQL clean and transform the data into ORC tables.
    \item \textbf{Analytics Layer:} Final HQL queries answer business questions
          on the cleaned tables.
\end{enumerate}

\subsection{Table Lineage}
\begin{center}
\texttt{emp\_raw} $\xrightarrow{\text{5 fixes + ROW\_NUMBER}}$
\texttt{emp\_clean} $\xrightarrow{\text{rn = 1}}$
\texttt{emp\_no\_dup} $\xrightarrow{\text{SPLIT + REGEXP}}$
\texttt{emp\_split\_title}
\end{center}

% Screenshot: Hadoop / HDFS running and the data directory
\begin{figure}[H]
\centering
\includegraphics[width=0.95\textwidth]{images/hdfs_upload.png}
\caption{HDFS upload verification -- \texttt{hdfs dfs -ls /data/employee/} confirming
         \texttt{EmployeeDataset.csv} is stored on the cluster.}
\label{fig:hdfs_upload}
\end{figure}

% ─────────────────────────────────────────────────────────────────────────────
\section{Task 1 -- Data Preprocessing (8 Marks)}
% ─────────────────────────────────────────────────────────────────────────────

\subsection{Raw Staging Table Creation}

All columns are declared as \texttt{STRING} so that dirty values (empty strings,
mixed casing) survive the ingestion without casting errors:

\begin{lstlisting}[caption={Creating the \texttt{emp\_raw} staging table and loading CSV from HDFS}]
CREATE TABLE emp_raw (
    job_id            STRING,
    job_title         STRING,
    experience_years  STRING,
    education_level   STRING,
    skills_count      STRING,
    industry          STRING,
    company_size      STRING,
    location          STRING,
    remote_work       STRING,
    call_log          STRING,
    salary            STRING,
    date_of_marriage  STRING
)
ROW FORMAT DELIMITED
FIELDS TERMINATED BY ','
STORED AS TEXTFILE
TBLPROPERTIES ("skip.header.line.count"="1");

LOAD DATA INPATH '/data/employee/EmployeeDataset.csv'
INTO TABLE emp_raw;
\end{lstlisting}

\begin{figure}[H]
\centering
\includegraphics[width=0.95\textwidth]{images/raw_table_load.png}
\caption{Hive output confirming \texttt{emp\_raw} table creation and CSV data load from HDFS.}
\label{fig:raw_load}
\end{figure}

\subsection{Five Preprocessing Problems Identified \& Fixed}

\subsubsection{Problem 1 -- Missing Values in \texttt{experience\_years}}

\textbf{Issue:} Two rows (job\_ids \textbf{1005} and \textbf{1006}) contain
empty \texttt{experience\_years} fields.  Loading them as-is would produce NULL
integers, silently skewing any average or filter.

\begin{lstlisting}[caption={Detecting missing \texttt{experience\_years}}]
SELECT job_id, job_title, experience_years, education_level
FROM emp_raw
WHERE experience_years IS NULL OR TRIM(experience_years) = '';
\end{lstlisting}

\begin{figure}[H]
\centering
\includegraphics[width=0.95\textwidth]{images/missing_experience.png}
\caption{Query result showing two rows with empty \texttt{experience\_years} -- job\_id 1005 (Business Analyst / PhD) and job\_id 1006 (Backend Developer / High School).}
\label{fig:missing_exp}
\end{figure}

\textbf{Fix:} Use \texttt{COALESCE(NULLIF(TRIM(...), ''), '0')} to substitute
empty strings with \texttt{0} before casting to \texttt{INT}.

\begin{lstlisting}[caption={Imputing missing experience with 0}]
CAST(COALESCE(NULLIF(TRIM(experience_years), ''), '0') AS INT)
    AS experience_years
\end{lstlisting}

\subsubsection{Problem 2 -- Duplicate \texttt{job\_id} Values}

\textbf{Issue:} \texttt{job\_id} should be unique, but IDs \textbf{1005},
\textbf{1006}, and \textbf{1009} each appear \emph{twice} with different roles
and data, suggesting data entry errors or system merges.

\begin{lstlisting}[caption={Identifying duplicate job\_ids}]
SELECT job_id, COUNT(*) AS cnt
FROM emp_raw
GROUP BY job_id
HAVING COUNT(*) > 1
ORDER BY cnt DESC;
\end{lstlisting}

\begin{figure}[H]
\centering
\includegraphics[width=0.95\textwidth]{images/duplicate_jobids.png}
\caption{Hive result listing duplicate \texttt{job\_id} values -- IDs 1005, 1006, and 1009 each appear twice.}
\label{fig:duplicates}
\end{figure}

\textbf{Fix:} Apply a \texttt{ROW\_NUMBER()} window function partitioned by
\texttt{job\_id} and keep only the first occurrence (\texttt{rn = 1}).

\begin{lstlisting}[caption={Deduplication using ROW\_NUMBER()}]
ROW_NUMBER() OVER (PARTITION BY job_id ORDER BY job_id) AS rn
-- Then filter: WHERE rn = 1
\end{lstlisting}

\subsubsection{Problem 3 -- Inconsistent Casing in \texttt{education\_level}}

\textbf{Issue:} The same qualification is recorded in multiple formats:
\texttt{high School}, \texttt{bachelor}, \texttt{PhD}, \texttt{Bachelor} --
causing group-by queries to split logically identical groups.

\begin{lstlisting}[caption={Revealing distinct raw education\_level values}]
SELECT DISTINCT education_level FROM emp_raw ORDER BY education_level;
\end{lstlisting}

\begin{figure}[H]
\centering
\includegraphics[width=0.95\textwidth]{images/inconsistent_casing.png}
\caption{Raw distinct values of \texttt{education\_level} -- note \texttt{high School}, \texttt{bachelor}, and \texttt{PhD} representing the same level with different casing.}
\label{fig:casing}
\end{figure}

\textbf{Fix:} Chain \texttt{LOWER()} then \texttt{INITCAP()} to canonicalize
every value (e.g.\ \texttt{"high School"} $\to$ \texttt{"High School"}).

\begin{lstlisting}[caption={Normalising education\_level casing}]
INITCAP(LOWER(TRIM(education_level))) AS education_level
\end{lstlisting}

\subsubsection{Problem 4 -- Outlier Salaries (Implausibly Low Values)}

\textbf{Issue:} Four rows contain salaries below \$1{,}000
(\textbf{50}, \textbf{99}, \textbf{100}, \textbf{102}), which are clearly
data entry errors -- likely missing digits (e.g.\ 50 instead of 50{,}000).

\begin{lstlisting}[caption={Detecting outlier salary values}]
SELECT job_id, job_title, salary
FROM emp_raw
WHERE CAST(salary AS BIGINT) < 1000;
\end{lstlisting}

\begin{figure}[H]
\centering
\includegraphics[width=0.95\textwidth]{images/outlier_salary.png}
\caption{Four rows flagged with implausibly low salary values (below \$1{,}000) -- likely data entry errors with missing digits.}
\label{fig:salary_outlier}
\end{figure}

\textbf{Fix:} Set these values to \texttt{NULL} so they are excluded from
aggregate calculations while the rows are retained for other columns.

\begin{lstlisting}[caption={Nullifying outlier salaries}]
CASE
    WHEN CAST(salary AS BIGINT) < 1000 THEN NULL
    ELSE CAST(salary AS BIGINT)
END AS salary
\end{lstlisting}

\subsubsection{Problem 5 -- \texttt{"Remote"} in the \texttt{location} Column}

\textbf{Issue:} Several rows use \texttt{"Remote"} as the value for
\texttt{location}, but \emph{Remote} describes a \textbf{work arrangement},
not a geographic country.  This contaminates any country-level aggregation
(e.g.\ Task 4 -- PhD count by country).

\begin{lstlisting}[caption={Identifying rows where location = "Remote"}]
SELECT job_id, location, remote_work
FROM emp_raw
WHERE LOWER(TRIM(location)) = 'remote';
\end{lstlisting}

\begin{figure}[H]
\centering
\includegraphics[width=0.95\textwidth]{images/remote_location.png}
\caption{Rows where \texttt{location} is set to \texttt{"Remote"} -- a work-mode value that should appear only in the \texttt{remote\_work} column.}
\label{fig:remote_loc}
\end{figure}

\textbf{Fix:} Replace \texttt{location = 'Remote'} with \texttt{'Unknown'} and
ensure \texttt{remote\_work} is forced to \texttt{'Yes'} for those rows.

\begin{lstlisting}[caption={Fixing the Remote location problem}]
CASE
    WHEN LOWER(TRIM(location)) = 'remote' THEN 'Unknown'
    ELSE TRIM(location)
END AS location,
CASE
    WHEN LOWER(TRIM(location)) = 'remote' THEN 'Yes'
    ELSE TRIM(remote_work)
END AS remote_work
\end{lstlisting}

\subsection{Final Cleaned Table}

All five fixes are combined in a single CTAS statement producing an ORC table
with Snappy compression:

\begin{lstlisting}[caption={Creating \texttt{emp\_no\_dup} -- the final cleaned table}]
CREATE TABLE emp_clean STORED AS ORC
TBLPROPERTIES ("orc.compress"="SNAPPY") AS
SELECT
    job_id,
    job_title,
    CAST(COALESCE(NULLIF(TRIM(experience_years),''),'0') AS INT) AS experience_years,
    INITCAP(LOWER(TRIM(education_level)))                        AS education_level,
    CAST(skills_count AS INT)                                    AS skills_count,
    industry, company_size,
    CASE WHEN LOWER(TRIM(location))='remote' THEN 'Unknown'
         ELSE TRIM(location) END                                 AS location,
    CASE WHEN LOWER(TRIM(location))='remote' THEN 'Yes'
         ELSE TRIM(remote_work) END                              AS remote_work,
    CAST(call_log AS INT)                                        AS call_log,
    CASE WHEN CAST(salary AS BIGINT) < 1000 THEN NULL
         ELSE CAST(salary AS BIGINT) END                         AS salary,
    TRIM(date_of_marriage)                                       AS date_of_marriage,
    ROW_NUMBER() OVER (PARTITION BY job_id ORDER BY job_id)      AS rn
FROM emp_raw;

-- De-duplicate: keep first occurrence per job_id
CREATE TABLE emp_no_dup STORED AS ORC
TBLPROPERTIES ("orc.compress"="SNAPPY") AS
SELECT job_id, job_title, experience_years, education_level,
       skills_count, industry, company_size, location,
       remote_work, call_log, salary, date_of_marriage
FROM emp_clean WHERE rn = 1;
\end{lstlisting}

\begin{figure}[H]
\centering
\includegraphics[width=0.95\textwidth]{images/cleaned_table.png}
\caption{Preview of the final \texttt{emp\_no\_dup} cleaned ORC table -- all five preprocessing
         problems resolved, data ready for analysis.}
\label{fig:cleaned_table}
\end{figure}

% ─────────────────────────────────────────────────────────────────────────────
\section{Task 2 -- Split Job Title into Two Columns (3 Marks)}
% ─────────────────────────────────────────────────────────────────────────────

The \texttt{job\_title} column contains multi-word strings
(e.g.\ \textit{``Machine Learning Engineer''}). We split on the \textbf{first
space} to produce two new columns:

\begin{table}[H]
\centering
\caption{Column Split Definition}
\begin{tabular}{@{}lll@{}}
\toprule
\textbf{New Column} & \textbf{Content} & \textbf{Example (from ``Machine Learning Engineer'')} \\
\midrule
\texttt{first\_title\_word} & First token & \texttt{Machine} \\
\texttt{remaining\_title}   & All other tokens & \texttt{Learning Engineer} \\
\bottomrule
\end{tabular}
\end{table}

\begin{lstlisting}[caption={Splitting \texttt{job\_title} using SPLIT and REGEXP\_REPLACE}]
CREATE TABLE emp_split_title STORED AS ORC
TBLPROPERTIES ("orc.compress"="SNAPPY") AS
SELECT
    job_id,
    job_title,
    -- First word (before the first space)
    SPLIT(TRIM(job_title), ' ')[0]
        AS first_title_word,
    -- Remainder (everything after the first space)
    REGEXP_REPLACE(TRIM(job_title),
        CONCAT('^', SPLIT(TRIM(job_title), ' ')[0], ' ?'), '')
        AS remaining_title,
    experience_years, education_level, skills_count,
    industry, company_size, location, remote_work,
    call_log, salary, date_of_marriage
FROM emp_no_dup;

SELECT job_id, job_title, first_title_word, remaining_title
FROM emp_split_title
LIMIT 20;
\end{lstlisting}

\begin{figure}[H]
\centering
\includegraphics[width=0.95\textwidth]{images/split_title.png}
\caption{Hive output of \texttt{emp\_split\_title} -- \texttt{job\_title} correctly decomposed
         into \texttt{first\_title\_word} and \texttt{remaining\_title} for all 20 sample rows.}
\label{fig:split_title}
\end{figure}

% ─────────────────────────────────────────────────────────────────────────────
\section{Task 3 -- Count PhD Holders with \textless 10 Years Experience by Country (2 Marks)}
% ─────────────────────────────────────────────────────────────────────────────

\begin{lstlisting}[caption={PhD employees with less than 10 years of experience, grouped by country}]
SELECT
    location            AS country,
    COUNT(*)            AS phd_under_10yrs_count
FROM emp_no_dup
WHERE LOWER(TRIM(education_level)) = 'phd'
  AND experience_years < 10
GROUP BY location
ORDER BY phd_under_10yrs_count DESC;
\end{lstlisting}

\begin{figure}[H]
\centering
\includegraphics[width=0.95\textwidth]{images/phd_experience_country.png}
\caption{Countries ranked by the number of PhD-qualified employees with fewer than
         10 years of professional experience.}
\label{fig:phd_country}
\end{figure}

% ─────────────────────────────────────────────────────────────────────────────
\section{Task 4 -- Skills Count in Ascending Order (1 Mark)}
% ─────────────────────────────────────────────────────────────────────────────

\begin{lstlisting}[caption={Employee records ordered by skills\_count ascending}]
-- Per-employee listing in ascending order of skill count
SELECT job_id, job_title, education_level, skills_count
FROM emp_no_dup
ORDER BY skills_count ASC;

-- Frequency distribution of skill counts
SELECT skills_count, COUNT(*) AS frequency
FROM emp_no_dup
GROUP BY skills_count
ORDER BY skills_count ASC;
\end{lstlisting}

\begin{figure}[H]
\centering
\includegraphics[width=0.95\textwidth]{images/skills_count_asc.png}
\caption{Employee records sorted by \texttt{skills\_count} in ascending order --
         ranging from 0 to 25 skills per employee.}
\label{fig:skills_asc}
\end{figure}

% ─────────────────────────────────────────────────────────────────────────────
\section{Task 5 -- Growth Rate of Call Log (3 Marks)}
% ─────────────────────────────────────────────────────────────────────────────

\subsection{Formula}

The row-to-row \textbf{Growth Rate} within each education group is:

\begin{equation}
\text{Growth Rate (\%)} =
\frac{\text{call\_log}_{\,n} - \text{call\_log}_{\,n-1}}
     {\text{call\_log}_{\,n-1}} \times 100
\label{eq:growth_rate}
\end{equation}

where rows are ordered by \texttt{job\_id} within each \texttt{education\_level} partition.

\subsection{HQL Implementation}

\begin{lstlisting}[caption={Computing row-to-row call\_log growth rate with the LAG window function}]
SELECT
    education_level,
    job_id,
    call_log AS current_call_log,
    LAG(call_log) OVER (
        PARTITION BY education_level
        ORDER BY job_id
    ) AS prev_call_log,
    ROUND(
        CASE
            WHEN LAG(call_log) OVER (
                     PARTITION BY education_level ORDER BY job_id
                 ) IS NULL
              OR LAG(call_log) OVER (
                     PARTITION BY education_level ORDER BY job_id
                 ) = 0
            THEN NULL
            ELSE ((call_log - LAG(call_log) OVER (
                              PARTITION BY education_level ORDER BY job_id
                          ))
                  / CAST(LAG(call_log) OVER (
                              PARTITION BY education_level ORDER BY job_id
                         ) AS DOUBLE)) * 100.0
        END
    , 2) AS growth_rate_pct
FROM emp_no_dup
WHERE LOWER(TRIM(education_level)) IN ('phd','bachelor','high school')
ORDER BY education_level, job_id;
\end{lstlisting}

\begin{figure}[H]
\centering
\includegraphics[width=0.95\textwidth]{images/growth_rate_detail.png}
\caption{Row-by-row call\_log growth rate for PhD, Bachelor, and High School employees
         -- computed using the \texttt{LAG()} window function partitioned by education level.}
\label{fig:growth_rate}
\end{figure}

\subsection{Average Growth Rate Summary}

\begin{lstlisting}[caption={Average growth rate per education group}]
SELECT education_level,
       ROUND(AVG(growth_rate_pct), 2) AS avg_growth_rate_pct
FROM (
    SELECT education_level,
           ROUND(CASE
               WHEN LAG(call_log) OVER (
                        PARTITION BY education_level ORDER BY job_id
                    ) IS NULL
                 OR LAG(call_log) OVER (
                        PARTITION BY education_level ORDER BY job_id
                    ) = 0
               THEN NULL
               ELSE ((call_log - LAG(call_log) OVER (
                                     PARTITION BY education_level ORDER BY job_id
                                 ))
                     / CAST(LAG(call_log) OVER (
                                 PARTITION BY education_level ORDER BY job_id
                            ) AS DOUBLE)) * 100.0
           END, 2) AS growth_rate_pct
    FROM emp_no_dup
    WHERE LOWER(TRIM(education_level)) IN ('phd','bachelor','high school')
) growth_sub
GROUP BY education_level
ORDER BY education_level;
\end{lstlisting}

\begin{figure}[H]
\centering
\includegraphics[width=0.95\textwidth]{images/growth_rate_summary.png}
\caption{Average call\_log growth rate summary per education level (PhD, Bachelor, High School).}
\label{fig:growth_summary}
\end{figure}

% ─────────────────────────────────────────────────────────────────────────────
\section{Task 6 -- Predict Future Call Log Values (3 Marks)}
% ─────────────────────────────────────────────────────────────────────────────

\subsection{Formula}

Future values are projected using the \textbf{Compound Growth Model}:

\begin{equation}
\text{Future Value}_{n} =
\text{Last Value} \times
\left(1 + \frac{\overline{g}}{100}\right)^{n}
\label{eq:future_value}
\end{equation}

where $\overline{g}$ is the average growth rate (\%) for the education group and
$n \in \{1, 2, 3\}$ represents future periods.

\subsection{HQL Implementation}

\begin{lstlisting}[caption={Predicting future call\_log values using a CTE and POWER()}]
WITH growth_stats AS (
    SELECT education_level, job_id, call_log,
           LAG(call_log) OVER (
               PARTITION BY education_level ORDER BY job_id
           ) AS prev_call_log
    FROM emp_no_dup
    WHERE LOWER(TRIM(education_level))
          IN ('phd', 'bachelor', 'high school')
),
growth_pct AS (
    SELECT education_level, job_id, call_log,
           CASE WHEN prev_call_log IS NULL OR prev_call_log = 0 THEN NULL
                ELSE ((call_log - prev_call_log)
                      / CAST(prev_call_log AS DOUBLE)) * 100.0
           END AS growth_rate_pct
    FROM growth_stats
),
summary AS (
    SELECT education_level,
           COUNT(*)                       AS n_records,
           ROUND(AVG(call_log), 4)        AS avg_call_log,
           MAX(call_log)                  AS last_call_log,
           ROUND(AVG(growth_rate_pct), 4) AS avg_growth_rate_pct
    FROM growth_pct
    GROUP BY education_level
)
SELECT
    education_level,
    n_records,
    avg_call_log,
    last_call_log,
    ROUND(avg_growth_rate_pct, 2)                                    AS avg_growth_pct,
    ROUND(last_call_log * POWER(1 + avg_growth_rate_pct/100.0, 1), 4) AS predicted_p1,
    ROUND(last_call_log * POWER(1 + avg_growth_rate_pct/100.0, 2), 4) AS predicted_p2,
    ROUND(last_call_log * POWER(1 + avg_growth_rate_pct/100.0, 3), 4) AS predicted_p3
FROM summary
ORDER BY education_level;
\end{lstlisting}

\begin{figure}[H]
\centering
\includegraphics[width=0.95\textwidth]{images/future_prediction.png}
\caption{Predicted future call\_log values for periods 1, 2, and 3 -- computed via compound
         growth from the historical average growth rate per education level.}
\label{fig:prediction}
\end{figure}

% ─────────────────────────────────────────────────────────────────────────────
\section{Script Files Summary}
% ─────────────────────────────────────────────────────────────────────────────

\begin{table}[H]
\centering
\caption{Project Scripts}
\begin{tabularx}{\textwidth}{@{}lX@{}}
\toprule
\textbf{Script} & \textbf{Purpose} \\
\midrule
\texttt{scripts/upload\_employee\_to\_hdfs.sh} &
    Starts HDFS, creates \texttt{/data/employee/} directory, and uploads
    \texttt{EmployeeDataset.csv} to HDFS. Verifies upload with \texttt{hdfs dfs -ls}. \\[4pt]
\texttt{scripts/employee\_analysis.hql} &
    Main Hive script.  Sections 0--7 cover table creation, all five
    preprocessing fixes, title split, and analytics queries. \\[4pt]
\texttt{scripts/run\_employee\_analysis.sh} &
    One-shot runner: calls the upload script then executes the HQL file
    via \texttt{hive -f}, piping output to a timestamped log file. \\
\bottomrule
\end{tabularx}
\end{table}

\textbf{To execute the complete pipeline:}
\begin{lstlisting}[language=bash, caption={Running the full analysis from the project root}]
bash scripts/run_employee_analysis.sh
\end{lstlisting}

% ─────────────────────────────────────────────────────────────────────────────
\section{Conclusion}
% ─────────────────────────────────────────────────────────────────────────────

This practical demonstrated the full power of the Hadoop/Hive ecosystem for
big data preprocessing and analysis.  Key takeaways:

\begin{itemize}
    \item \textbf{Five preprocessing problems} were systematically identified
          and resolved using HQL built-in functions
          (\texttt{COALESCE}, \texttt{INITCAP}, \texttt{ROW\_NUMBER}, \texttt{CASE}).
    \item \textbf{Window functions} (\texttt{LAG}, \texttt{ROW\_NUMBER})
          enable row-to-row comparisons without leaving the Hive SQL paradigm.
    \item \textbf{String manipulation} (\texttt{SPLIT}, \texttt{REGEXP\_REPLACE})
          allows flexible column decomposition (Task 2 -- title split).
    \item \textbf{CTEs} (\texttt{WITH} clauses) improve query readability for
          multi-step analytical workflows (Task 6 -- future prediction).
    \item \textbf{ORC + Snappy} storage reduces disk footprint and accelerates
          subsequent reads compared to plain text files.
\end{itemize}

The compound growth model used for prediction is straightforward and interpretable,
making it well-suited for a business analytics context.  For production use on
larger, time-indexed datasets, a regression or Holt-Winters time-series model
would provide more accurate forecasts.

\end{document}
